\section{Message Deletion in the Wild}\label{sec:attacks}
Equipped with our threat model, we now turn to the motivating question behind this work: Do real-world applications safely delete messages? We explore this question through the lens of two case-study applications, WhatsApp and Signal. Our security analysis reveals an unsatisfactory landscape: For both applications, we are able to recover deleted messages.

\paragraph{Attacks overview}

\paragraph{Methodology and experimental setup}

\paragraph{Architecture of WhatsApp and Signal}

\paragraph{Message deletion features in WhatsApp and Signal}

\paragraph{Recovery of application state}

\paragraph{FTS index}

\subsection{WhatsApp}

\paragraph{Architecture.} \andres{Start by saying we focus on Android but believe our attacks generalize. Or maybe say this in a methodology paragraph.}

\paragraph{Message deletion features.}

\andres{Here, explain recovery of application state.}
\andres{First, an overview of general deletion. What table changes. Then what are issues with this.}

\paragraph{Vulnerability \#1: write-ahead log.}

\paragraph{Vulnerability \#2: FTS index.}

\paragraph{Vulnerability \#3: view-once media.}

\paragraph{Vulnerability \#4: attachment deduplication.}

\paragraph{Vulnerability \#5: message reactions.}

\paragraph{Vulnerability \#6: pinned messages.}

\subsection{Signal}

\paragraph{Architecture.}

\paragraph{Message deletion features.}

\paragraph{Vulnerability \#1: freelist pages}

\paragraph{Vulnerability \#2: attachment deduplication.}



